\section{A nonlinear first-order family near the characteristic cutoff}
\label{app:inverse-bernoulli}

The bounded-root result of Appendix~\ref{app:bounded-root-mca} requires
characteristic larger than a suitable constant times the length.
For the inverse Bernoulli family below, a direct classification also
handles characteristic just above the message degree. The extra
Frobenius constants available in that range do not produce large lists.

\begin{theorem}[Inverse Bernoulli solution count]
\label{thm:inverse-bernoulli}
Let $s\ge2$, $D\ge1$, and let $E$ be algebraically closed of
characteristic zero or $p>\max\{D,s+1\}$. For every nonzero polynomial
$A(X)$, the equation
\[
 P^{s-1}P'=A,\qquad \deg P\le D,
\]
has at most $2s$ polynomial solutions, all regular. In characteristic
zero, or when $p>sD$, it has at most $s$ solutions. The bound $2s$
is attained for suitable degrees and primes with $p\equiv1\pmod s$.
\end{theorem}
\begin{proof}
Every solution has a positive degree $d<p$ (with the obvious
characteristic-zero interpretation), and
$\deg A=sd-1$. If its leading coefficient is $a$, the leading
coefficient of $A$ is $da^s$. Thus all solutions have the same degree
and their leading coefficients differ by $s$th roots of unity.
Normalize one leading coefficient to one. It suffices to count
monic solutions of the resulting equation and multiply by $s$.
For two such solutions $P,Q$,
\[
 (P^s-Q^s)'=0.
\]
In characteristic zero their difference of powers is constant.
Factoring it as $\prod_{\zeta^s=1}(P-\zeta Q)$ shows that it cannot
be a nonzero constant: a factor with $\zeta\ne1$ has degree $d$.
If it is zero, monicity gives $P=Q$. The same argument applies when
$p>sD$. Otherwise, in positive characteristic, a distinct pair has
\[
 P^s-Q^s=C(X)^p,\qquad 1\le\deg C\le s-1.
\]

\paragraph{At least three factors.}
Suppose $s\ge3$. Cancel the monic gcd $G$ of $P,Q$ and put
$U=P/G$, $V=Q/G$, both monic of the same degree $h\ge1$.
The rational map $U/V$ has degree $h<p$ and is separable.
The preimages of the $s$ distinct $s$th roots of unity all lie
among the roots of $C$ and infinity. These are at most $s$ points;
the $s$ fibers are disjoint and nonempty. Consequently each fiber
has exactly one point, of ramification index $h$.

For completeness, an elementary degree count forces $h=1$.
Make a generic fractional linear change of the $X$ coordinate so
that all those fiber points are finite and the new infinity maps
to none of the $s$ values, nor to zero or infinity. The new numerator
and denominator have degree $h$. At each of the $s$ points their
derivative numerator $U'V-UV'$ vanishes to order $h-1$, since
$h<p$. Its degree is at most $2h-2$. Hence
$s(h-1)\le2h-2$, which for $s\ge3$ gives $h=1$.

Return to the original coordinate, and write $U=X-u$, $V=X-v$,
with $u\ne v$. The polynomial $S=U^s-V^s$ has $s-1$ distinct
roots, none equal to $u$ or $v$. These exhaust the roots of $C$,
which therefore has degree $s-1$ and is squarefree. Comparing
valuations in $G^sS=C^p$ proves
\[
 p=sg+1,\qquad G=S_0^g,\qquad d=(s-1)g+1,
\]
where $S_0$ is the monic normalization of $S$ and $g=(p-1)/s$.
There are no other roots of $G$. The strict assumption $p>s+1$
ensures $g\ge2$.

Direct differentiation, using $sg=-1$ in the field, gives
\[
 A=\text{nonzero constant}\cdot (UV)^{s-1}S_0^{p-2}.
\]
At a root of multiplicity $m$ in any other solution $T$, the
multiplicity in $T^{s-1}T'$ is $sm-1$, since $m<p$.
Thus $T$ can choose roots only from $S_0$, with multiplicity $g$,
and from $u,v$, with multiplicity one. If it chooses $t$ of the
former and $k$ of the latter, its degree requires
\[
 tg+k=(s-1)g+1,\qquad 0\le t\le s-1,\quad0\le k\le2.
\]
For $g\ge2$ this forces $t=s-1$ and $k=1$, so $T=P$ or $Q$.
There are at most two monic solutions. Conversely the displayed
$U,V,S_0$ construction supplies two, and their $s$ scalar multiples
attain the bound whenever $p=sg+1$ with $g\ge2$.

\paragraph{Two factors.}
For $s=2$, a distinct monic pair satisfies
$(P-Q)(P+Q)=c(X-a)^p$. Since $P+Q$ has degree $d$ and leading
coefficient two, its two factors give
\[
 P=(X-a)^d+v(X-a)^{p-d},\qquad
 Q=(X-a)^d-v(X-a)^{p-d},\qquad v\ne0.
\]
In particular $d>p/2$. If $p-d\le d-2$, the coefficient of
$X^{d-1}$ in $P$ fixes $a$, and then $v$ is fixed. If $p-d=d-1$,
then $p=2d-1\ge5$ and $a$ is a root of $P$ of multiplicity
$d-1>d/2$, so it is again unique. Every monic solution has at most
one distinct partner. The displayed pair and its negatives attain
four solutions. Finally, the separant $P^{s-1}$ is nonzero for
every solution of a nonzero-$A$ equation, proving regularity.
\end{proof}

\paragraph{The strict characteristic condition has a real exception.}
For $s=p-1$ and $D=p-1$, the polynomials
\[
 P_a=\frac{X^p-X}{X-a},\qquad a\in\mathbb F_p,
 \qquad P_a^{p-2}P_a'=-(X^p-X)^{p-2},
\]
and all their nonzero $\mathbb F_p$ scalar multiples give $p(p-1)$
solutions. Here the jet degree $s$ grows with $p$. This does not
contradict Theorem~\ref{thm:inverse-bernoulli}, which assumes $p>s+1$.

\begin{corollary}[Linear full MCA for inverse Bernoulli families]
\label{cor:inverse-bernoulli-mca}
Under the theorem's hypotheses, for nonzero $A$ the solution family
witnesses at most $2sn$ full-support bad labels on any received line.
For $A=0$, the corresponding count at agreement threshold $a\ge2$
is at most $\binom n2/(a-1)$. In particular both counts are $O(n)$
at fixed $s$ and fixed positive agreement fraction.
\end{corollary}
\begin{proof}
A fixed candidate has at most $n$ bad labels: each bad full support
contains an agreeing coordinate with nonzero received direction,
which determines the label; otherwise $(P,0)$ explains that support.
Apply Theorem~\ref{thm:inverse-bernoulli} when $A\ne0$.

When $A=0$, the degree guard makes every solution constant.
Each coordinate gives an affine line $c=f(x)+zg(x)$ in the
$(z,c)$ plane. A bad support of size at least $a$ contains at least
two distinct line equations, since otherwise constant polynomials
$F,G$ explain the whole support. Among its coordinates there are
at least $a-1$ pairs with distinct line equations. Each such pair
intersects at at most one point. Select one candidate per bad label
and count pairs to obtain the stated bound.
\end{proof}

\paragraph{Checks and scope.}
The exact standard-library checker in
\path{research/inverse_bernoulli/} exhausts $164{,}803$ monic
polynomials across thirteen fixtures, checks six sharpness examples,
and verifies the exceptional family at $p=3,5,7,11$.
The argument is specific to inverse Bernoulli equations. It does
not settle the first-order conjecture or give a quadratic MCA
construction at any higher order.
